\section*{第 \ref{ch:g12:internal-heat} 章——地球内部热量}

\begin{solution}{exo:g12:internal-heat:1}
(a) \omterm{def:g12:internal-heat:heat-flow}{热流}；(b) \omterm{def:g12:internal-heat:heat-production}{放射性生热}；(c) \omterm{def:g12:internal-heat:secular-cooling}{长期冷却}；
(d) \omterm{def:g12:internal-heat:convection}{热对流}；(e) \omterm{def:g12:internal-heat:geothermal}{地热能}。
\end{solution}

\begin{solution}{exo:g12:internal-heat:2}
(1) \qty{0.0916}{W/m^2}。(2) \qty{4.67e13}{W}。(3) \qty{0.029}{K/m}：
一摄氏度的温差即为一开尔文的温差。
\end{solution}

\begin{solution}{exo:g12:internal-heat:3}
钍-232、铀-238、钾-40 和铀-235。KamLAND 能探测到
仅钍-232 和铀-238 的地中微子。
\end{solution}

\begin{solution}{exo:g12:internal-heat:4}
$490/\sqrt{1} = \qty{490}{mW/m^2}$；$490/2 = \qty{245}{mW/m^2}$；
$490/5 = \qty{98}{mW/m^2}$；$490/7 = \qty{70}{mW/m^2}$，
约为大陆的平均值。
\end{solution}

\begin{solution}{exo:g12:internal-heat:5}
近地表处，地面温度随昼夜和季节变化；
只有在几米以下，温度才稳定在全年的平均温度，
直到更深处，才显现来自下方的热量。
钻探会搅动钻孔中的岩石和水，改变它们的
温度，因此需等待钻孔恢复到
岩石的温度。
\end{solution}

\begin{solution}{exo:g12:internal-heat:6}
(1) $(265 - 10)/9.1 \approx \qty{28.0}{\celsius}$ 每千米，即
\qty{0.0280}{K/m}。(2) $q = 2.72 \times 0.0280 \approx \qty{0.0762}{W/m^2}$，
约 \qty{76}{mW/m^2}。(3) 略高于大陆的平均值，
\qty{70.9}{mW/m^2}，约高出 \qty{7}{\percent}，为典型的大陆\omterm{def:g12:internal-heat:heat-flow}{热流}值。
\end{solution}

\begin{solution}{exo:g12:internal-heat:7}
大陆：$0.0709 \times \num{2.073e14} \approx \qty{1.47e13}{W} =
\qty{14.7}{TW}$。海洋：$0.1054 \times \num{3.028e14} \approx
\qty{3.19e13}{W} = \qty{31.9}{TW}$。总计约 \qty{46.6}{TW}，其中
$31.9/46.6 \approx \qty{68}{\percent}$，约三分之二，
经洋底散失。
\end{solution}

\begin{solution}{exo:g12:internal-heat:8}
(1) $\num{1.88e-6} \times 8000 \approx \qty{0.0150}{W/m^2} = \qty{15.0}{mW/m^2}$。
(2) $15.0/70.9 \approx \qty{21}{\percent}$，约为\omterm{def:g12:internal-heat:heat-flow}{热流}的五分之一。
\end{solution}

\begin{solution}{exo:g12:internal-heat:9}
$\sqrt{t} = 490/70.9 \approx 6.91$，因此 $t \approx 47.8$ 百万年：洋底
年龄超过约 48 百万年时，散热少于
平均大陆的散热。
\end{solution}

\begin{solution}{exo:g12:internal-heat:10}
铀-235：$2^{1.4/0.704} = 2^{1.99} \approx 4$ 倍。钾-40：
$2^{2.5/1.248} = 2^{2.00} \approx 4$ 倍。两者均相当于回溯两个半衰期。
\end{solution}

\begin{solution}{exo:g12:internal-heat:11}
(1) $10^2/10^{-6} = \qty{1e8}{s}$，约 3 年。(2) $(10^5)^2/10^{-6} =
\qty{1e16}{s}$，约 3.2 亿年。(3) $(\num{2.891e6})^2/10^{-6}
\approx \qty{8.4e18}{s}$，约 2650 亿年。时间随距离的平方增长，
厚度增至十倍，时间则增至一百倍。
\end{solution}

\begin{solution}{exo:g12:internal-heat:12}
各乘以 $2^{3/T}$：钍-232，$10.4 \times 2^{3/14.0} \approx
\qty{12.1}{TW}$；铀-238，$9.6 \times 2^{3/4.468} \approx \qty{15.3}{TW}$；
钾-40，$4.0 \times 2^{3/1.248} \approx \qty{21.2}{TW}$；铀-235，
$0.42 \times 2^{3/0.704} \approx \qty{8.1}{TW}$。总计约 \qty{56.7}{TW}，
是今日 \qty{24.4}{TW} 的两倍以上。两种短半衰期同位素
占其中 $(21.2 + 8.1)/56.7 \approx \qty{52}{\percent}$，相比之下，
今日为 $4.42/24.4 \approx \qty{18}{\percent}$。
\end{solution}

\begin{solution}{exo:g12:internal-heat:13}
功率等于每秒的质量、比热与
温升的乘积：$\num{3e8} = 1640 \times 4180 \times \Delta T$，因此 $\Delta T =
\num{3e8}/(1640 \times 4180) \approx \qty{44}{K}$。大陆平均每平方米散失
\qty{0.0709}{W}，因此 \qty{300}{MW} 相当于
面积 $\num{3e8}/0.0709 \approx \qty{4.2e9}{m^2}$ 上的自然散热量，约 \qty{4200}{km^2}：
电站所收集的是地球在数千
平方千米范围内自然释放的热量。
\end{solution}

\begin{solution}{exo:g12:internal-heat:14}
(1) 1 亿年比 2600 亿年短 $\num{2.6e11}/\num{1e8} = 2600$ 倍；
2 亿年则短 1300 倍。
(2) 对流借运动的岩石输送热量：穿越
地幔所需时间等于其厚度除以速度，即 $\qty{2891}{km}$ 按
每年几厘米计，约一亿年。传导没有这样的
载体，其时间随距离的平方增长，在数千
千米尺度上会变得极其漫长。
\end{solution}

\begin{solution}{exo:g12:internal-heat:15}
(1) $10/22.3 \approx \qty{45}{\percent}$；其余为
地幔自身冷却，即失去其储存的热量。(2) 边界位于半径 $6371 - 2891 =
\qty{3480}{km}$；其面积为 $4\pi (\num{3.48e6})^2 \approx
\qty{1.52e14}{m^2}$，$\num{1e13}/\num{1.52e14} \approx
\qty{0.066}{W/m^2}$，约 \qty{66}{mW/m^2}，约占
地表平均值的 \qty{72}{\percent}。(3) 热量通过传导穿过边界，进入
地幔最底层；岩石受热后，比其上方的
地幔更轻，可以上升并携带热量向上输送：这被认为是
地幔柱在深部开始形成的方式。
\end{solution}

\begin{solution}{pb:g12:internal-heat:1}
\textbf{1.} $4\pi (\num{6.371e6})^2 \approx \qty{5.101e14}{m^2}$，
而 $2.073 + 3.028 = 5.101$，成立。

\textbf{2.} $3.028/5.101 \approx \qty{59}{\percent}$，约为五分之三。

\textbf{3.} 洋底很年轻：它在洋中脊处被加热，
随远离洋脊而冷却，年轻的\omterm{def:g11:lithosphere-asthenosphere:lithosphere}{岩石圈}散失了大量热量。

\textbf{4.} $490/\sqrt{4} = \qty{245}{mW/m^2}$。

\textbf{5.} $2.72 \times 0.026 \approx \qty{0.0707}{W/m^2} =
\qty{70.7}{mW/m^2}$，略低于平均值，实际上与之基本相等。

\textbf{6.} $0.0709 \times \num{2.073e14} \approx \qty{14.7}{TW}$。

\textbf{7.} $0.1054 \times \num{3.028e14} \approx \qty{31.9}{TW}$。

\textbf{8.} $14.7 + 31.9 \approx \qty{46.6}{TW}$，与已发表估计的
\qty{47}{TW} 相符（本章全球平均值为 \qty{46.7}{TW}，
即下文所用的值）。

\textbf{9.} $\num{4.66e13}/\num{5.101e14} \approx \qty{0.0914}{W/m^2}$，
约为 \qty{91}{mW/m^2}。

\textbf{10.} $0.0243 \times \num{2.073e14} \approx \qty{5.0}{TW}$，约为
通过大陆散失的 \qty{14.7}{TW} 的三分之一：\omterm{def:g12:continental-crust:continental-crust}{大陆地壳}的放射性
是大陆\omterm{def:g12:internal-heat:heat-flow}{热流}的很大一部分。

\textbf{11.} 地球中每千克铀对应含有 \qty{3.9}{kg} 钍，
这些钍释放 $3.9 \times 26.3 \approx \qty{102.6}{\micro W}$，
而铀-238 释放 \qty{94.4}{\micro W}。因此铀-238 占
$94.4/(94.4 + 102.6) \approx 0.479$ 的 \qty{20.0}{TW}，约为
\qty{9.6}{TW}。铀-235：$9.59 \times 4.10/94.4 \approx \qty{0.42}{TW}$。

\textbf{12.} $20.0 + 0.42 + 4 \approx \qty{24.4}{TW}$。

\textbf{13.} 花岗岩及\omterm{def:g12:continental-crust:continental-crust}{大陆地壳}中的其他岩石，
其放射性元素的富集程度约为地幔超\omterm{def:g11:igneous-rocks:felsic-mafic}{镁铁质}岩石的二十倍。

\textbf{14.} 更大：当时每种同位素都更丰富，每回溯一个半衰期，其量便多一倍；
30 亿年前的总量超过
今天的两倍。

\textbf{15.} $46.7 - 24.4 = \qty{22.3}{TW}$。

\textbf{16.} 来自地球形成以来储存的热量，即
地幔与地核的\omterm{def:g12:internal-heat:secular-cooling}{长期冷却}。

\textbf{17.} $22.3/46.7 \approx \qty{48}{\percent}$。

\textbf{18.} 从 $11/46.7 \approx \qty{24}{\percent}$ 到 $38/46.7 \approx
\qty{81}{\percent}$，地球中微子极大地收窄了这一范围。

\textbf{19.} \omterm{def:g12:internal-heat:convection}{地幔对流}输送了大部分热量；
只有最后数十千米由传导穿过。

\textbf{20.} $24.4/46.7 \approx 0.52$：放射性提供今日地球热损失的
\textbf{约一半}，
其余一半由行星的缓慢冷却提供。
\end{solution}
